\chapter{Introduction} \label{sec:Introduction}

\section{Context And Motivation} \label{sec:ContextAndMotivation}

The video game industry, first emerging in the 1970s and 1980s, is now the fastest
growing sector of the entertainment industry, and is expected to surpass a value of
US\$ 521 billion by 2027. Nearly 40\% of the population spends at least some of their
leisure time gaming. \cite{goh_unravelling_2023} With the rise of video games, a rise of interest in video game
development has followed. One of the primary aspects of games that developers
often work with is writing narratives and characters. Linear narratives have been
a practiced medium for millenia now, and are therefore well-understood and easy
to implement. However, branching or emergent narratives are a much newer and
less understood area. Game developers writing these narratives often have trouble
managing the exponential complexity of such a narrative, and the games that have
done it successfully are backed by large studios with hundreds of employees.
Over the years, tools have been both proposed and developed for this purpose, and
although many have promising features, few have had a significant impact on the wider
game development scene. This paper proposes a new solution that draws from existing
ideas, hoping to offer a framework for writing branching and emergent narratives that
can be widely applied to existing systems.

\section{Research Questions} \label{sec:ResearchQuestions}

\begin{itemize}
    \item What elements make up a successful character-focused narrative DSL?
    \item Is a structured relationship graph a useful way to model world state for
    writing autonomous characters?
    \item Can the fundamentals of Praxis and Comme il Faut be created
    in a modernized, type-safe scripting language?
\end{itemize}

\section{Hypothesis} \label{sec:Hypothesis}

A modern domain-specific language can be created that is intuitive for authors to
use, while also allowing for extensive flexibility and the ability to create emergent
narratives by putting characters at the center of attention.

\section{Contributions} \label{sec:IntroContributions}

This paper aims to contribute an overview of what technologies and principles outline
a successful character-narrative DSL and interpret those in a modern context. The
implementation developed to address these points, Praxsmith, will include:

\begin{itemize}
    \item A flexible game engine built on top of a Rust backend and a web app frontend.
    \item An actor-centric DSL inspired by Praxis, as well as many other influences
    outlined in \Cref{sec:BackgroundAndLiteratureReview}.
    \item Modern tooling, such as build-time type-checking to prevent runtime errors
    and minimize logic bugs.
    \item An authoring environment built into the web app.
    \item An evaluation of usability and capability, comprising of author evaluation
    plus a limited user study.
\end{itemize}
